\documentclass[11pt]{article}
\usepackage[margin=1in]{geometry}
\usepackage[T1]{fontenc}
\usepackage{lmodern,microtype,amsmath,amssymb,amsthm,mathtools,mathrsfs}
\usepackage{enumitem,booktabs,array,needspace,etoolbox}
\usepackage[colorlinks=true,linkcolor=blue,citecolor=blue,urlcolor=blue]{hyperref}
\usepackage{fancyhdr}
\pagestyle{fancy}\fancyhf{}
\fancyhead[L]{\small Focused Diophantine document}
\fancyhead[R]{\small Research draft}
\fancyfoot[C]{\thepage}
\setlength{\headheight}{14pt}
\setlength{\emergencystretch}{2em}
\newtheorem{theorem}{Theorem}[section]
\newtheorem{lemma}[theorem]{Lemma}
\newtheorem{proposition}[theorem]{Proposition}
\newtheorem{corollary}[theorem]{Corollary}
\theoremstyle{definition}
\newtheorem{definition}[theorem]{Definition}
\newtheorem{remark}[theorem]{Remark}
\newcommand{\Q}{\mathbb Q}\newcommand{\Z}{\mathbb Z}
\newcommand{\C}{\mathbb C}
\newcommand{\N}{\mathbb N}\newcommand{\F}{\mathbb F}
\newcommand{\Loc}{\operatorname{Loc}}
\newcommand{\Ctl}{\operatorname{Ctl}}
\newcommand{\Mult}{\operatorname{Mult}}
\newcommand{\Int}{\operatorname{Int}}
\newcommand{\Norm}{\operatorname{N}}
\newcommand{\Tr}{\operatorname{Tr}}
\newcommand{\Supp}{\operatorname{Supp}}
\newcommand{\Fin}{\operatorname{Fin}}
\newcommand{\Formal}{\operatorname{Form}}
\newcommand{\code}{\operatorname{code}}
\hypersetup{pdftitle={Focused Diophantine document: an explicit definition of the integers},pdfauthor={Research note prepared for James Freitag}}
\title{Focused Diophantine document\\\large An explicit definition of the integers in the rationals}
\author{Research note prepared for James Freitag}
\date{9 October 2026}
\begin{document}
\maketitle
\begin{abstract}
Using the arithmetic results of the cited OpenAI manuscript, we construct
an existential definition of the integers in the rationals. A finite
parity certificate defines multiplication on an elliptic integer model;
height bounds and a bounded divisibility certificate then exclude
denominators. Explicit sequence encodings give one specified polynomial
of degree four with \(17\,030\,326\,787\,289\) rational witnesses. The main text
gives the construction and proof; the appendices verify its numerical data.
\end{abstract}

\section{The definition and the arithmetic data}\label{sec:inputs}
\label{sec:dependency-status}
\begin{theorem}\label{thm:main}
Using the cited arithmetic results, the integer polynomial
\(P_{\rm final}\) specified in Section~\ref{sec:completed-definition}
has degree four and satisfies
\[
 x\in\Z\quad\Longleftrightarrow\quad
 \exists\mathbf w\in\Q^{17\,030\,326\,787\,289}\quad
                         P_{\rm final}(x,\mathbf w)=0.
\]
\end{theorem}

The construction develops the notes supplied by Gabe Conant
\cite{Conant,ConantArity}. We first define multiplication on an elliptic
model, then prove the integrality formula \eqref{eq:integer-formula}
and compile it into \(P_{\rm final}\). All quantified coordinates
range over \(\Q\); integer indices, primes and heights occur only
in the proof and in intermediate specifications that are eliminated.
The argument and exact computational checks remain a research draft,
without independent expert review or formal verification of the whole proof.

Fix
\[
 F=\Q(\sqrt2),\quad d=75-53\sqrt2,\quad E:y^2=x^3-d^2x,
 \qquad \sigma(\sqrt2)=-\sqrt2,
\]
\[
 P=\left(\frac{1641}{8}-145\sqrt2,
          -\frac{15465}{4}+\frac{87483}{32}\sqrt2\right),
 \qquad G=2E(F).
\]
A point code \((\tau,x_0,x_1,y_0,y_1)\) is either the zero
tuple, representing \(O\), or has \(\tau=1\) and represents
\((x_0+x_1\sqrt2,y_0+y_1\sqrt2)\). Write \(h\) for absolute
logarithmic Weil height, normalize \(v_p(p)=1\), and include zero
in \(\N\). Put
\[
 S=S_*:=\{p:p\le76176\},\quad K=48,\quad
 \mathcal C=\{\infty,0,1,2401,-343/32\},
\]
with contact forms \(L_\infty=V\) and
\[
 (L_1,L_2,L_3,L_4)=(U,U-V,U-2401V,32U+343V).
\]
Let \(H_*,L_*=18H_*+12\) be the explicit integers in
\eqref{eq:optimized-H}. We need the following four facts.

\begin{enumerate}[label=\textbf{(I\arabic*)},leftmargin=*,itemsep=4pt]
\item \emph{Elliptic data.} \(G=\Z P\), and
\begin{equation}\label{eq:canonical}
                    h(x(nP))\ge n^2-4\qquad(n\ne0).
\end{equation}
The rank assertion is \cite[Proposition 2.1]{H10}; Appendix
\ref{sec:numerical-elliptic} verifies the generator and height bound.
Set \(z=-x/y\), \(f=1/x\), with \(z(O)=f(O)=0\), and
let \(\ell_K\) be the truncated formal logarithm in
\eqref{eq:closed-log}. At a place over \(q\notin S_*\), for
formal multiples with denominator index \(n'\ne0\),
\begin{align}
 v\!\left(\frac{z(nP)}{z(n'P)}\right)&=v_q(n/n'),\label{eq:formal-ratio}\\
 v\!\left(\frac{\ell_K(z(nP))}{\ell_K(z(n'P))}-\frac n{n'}\right)
 &\ge K\quad\text{if }v_q(n/n')\ge0.\label{eq:formal-log}
\end{align}
These also hold on the conjugate curve \cite[Proposition 2.2 and
Section 5]{H10}: the formal logarithm preserves valuation, and its
tail has order at least \(v(z)+K\), since \(q>K+1\).

\item \emph{Parity.} The polynomial predicate \(\Phi\) in
Appendix~\ref{sec:parity} satisfies
\(D\subseteq\Phi(\Q)\subseteq D_{S_*}\), where \(D\) consists
of rationals with square reduced denominator, and \(D_{S_*}\)
allows odd denominator valuations only in \(S_*\).
It uses \cite[Proposition 2.8; Lemmas 3.4 and 3.6, with the local
conditions of Lemma 3.2]{H10}. The latter lemmas give one squarefree
parameter and points for both twists, as required in
\cite[Arithmetic Input 2.3]{Conant}.

\item \emph{Contact height.} For \(s=(u:v)\notin\mathcal C\),
the contact depth at \(b\) is
\(v_q(L_b(u,v))-\min(v_q(u),v_q(v))\).
Let \(M_S(s)\) be the product of primes outside \(S\) at which
some depth is odd. Appendix~\ref{sec:height-certificate} specializes
\cite[Theorem 2.4]{H10} to give
\begin{equation}\label{eq:height}
                      h(s)\le H_*M_S(s)^{24}.
\end{equation}
Appendix~\ref{sec:contacts} identifies the contacts from the
published quotient models \cite{PadurariuSaia}.

\item \emph{Ordinary contact witnesses.} Every integer \(A\)
admits integers \(u_1,u_2,u_3,v\), with \(v\ne0\) and
\(u_1+u_2+u_3=Av\), whose thirteen raw contacts
\begin{equation}\label{eq:thirteen}
              v,\quad L_b(u_i,v)\quad(1\le i\le3,\ 1\le b\le4)
\end{equation}
are \(\varepsilon kr\). Here \(\varepsilon=\pm1\), the
multiplier \(k\) belongs to the fixed finite catalogue
\eqref{eq:cyclic-multiplier}, and \(r\) is a positive split prime
exceeding any prescribed fixed lower bound. At each place over \(r\)
there are nonzero integers \(h,j\), independent of \(a\in\Z\),
such that \(hP,haP,jP^\sigma,(h-4)aP^\sigma\) are formal and
\begin{equation}\label{eq:ordinary-tests}
 \frac{\ell_K(z(haP))}{\ell_K(z(hP))}-a\in r^K\mathcal O_v,
 \qquad
 \frac{z((h-4)aP^\sigma)}{z(jP^\sigma)}\in r^K\mathcal O_v.
\end{equation}
The denominators are nonzero. This is
\cite[Corollaries 2.6 and 6.2, using Lemma 6.1]{H10}, with the
catalogue specialization proved in Appendix~\ref{sec:contacts}.
\end{enumerate}

All references to \cite{H10} use its 6 October 2026 version.
The appendices verify every new numerical specialization; the source
theorems themselves are used as cited results.

\section{Multiplication on the elliptic model}
\label{sec:simultaneous}\label{sec:multiplication}

The batching argument is adapted from \cite{ConantArity}. A
nonidentity point is formal at a good place exactly when \(v(f)>0\).

\subsection{One parity predicate for a list of inequalities}

\begin{lemma}[Simultaneous parity certificate]\label{lem:batch}
Let \(r\in\Q^\times\), \(X_1,\ldots,X_N\in\Q\), and
\(L_i\ge1\) be integers. Choose a power of two \(M>N\), and set
\[
 W=1+\sum_{i=1}^N r^i\left(\frac{X_i}{r^{L_i}}\right)^M.
\]
If \(\Phi(rW^2)\) holds and \(q\notin S_*\) has odd
\(v_q(r)\), then \(v_q(r)>0\) and \(v_q(X_i)\ge L_i\)
for every \(i\). Conversely, if \(r\) is a positive prime and
\(v_r(X_i)\ge L_i\) for all \(i\), then \(rW^2\in D\),
where \(D\) is the set of rationals with square reduced denominator.
\end{lemma}
\begin{proof}
Write \(R=v_q(r)\). The nonzero summands have pairwise distinct
valuations modulo \(M\), namely \(0,R,\ldots,NR\), because
\(R\) is odd. Their minimum is attained once. Consequently
\[
 W\ne0,\qquad v_q(W)=\min\bigl(0,\ iR+M(v_q(X_i)-L_iR)\bigr)\le0,
\]
where zero summands are omitted. If \(R<0\), the valuation
\(R+2v_q(W)\) is negative and odd, contradicting the parity
property of \(\Phi\). Thus \(R>0\); the same property gives
\(v_q(W)\ge-R/2\). For every nonzero \(X_i\),
\[
 v_q(X_i)\ge\left(L_i-\frac{i+1/2}{M}\right)R>L_i-1.
\]
The parenthesis is positive, \(R\ge1\), and \(M\ge N+1\).
Integral valuations give the conclusion; zero inputs satisfy it as
well. Conversely each nonconstant summand has positive valuation at
the prime \(r\), so \(W\) is an \(r\)-adic unit. At every
other prime the valuation of \(rW^2\) is even.
\end{proof}

For a field expression \(D=A+B\sqrt2\) required to have order
at least \(j\) at a selected place, use the two rational tests
\begin{equation}\label{eq:batch-field-tests}
                   (s^2-2)B,\qquad A+Bs
                   \quad\text{with orders }j,j.
\end{equation}
Also test \(s^2-2\) at order one. For an odd prime selected by
Lemma~\ref{lem:batch}, this supplies a root \(\theta\in\Q_q\)
with \(v_q(\theta-s)=v_q(s^2-2)\): both \(s\) and
\(\theta+s\) are units. Hence
\(v_q(B(\theta-s))\ge j\), and \(v_q(A+B\theta)\ge j\).
There is no assumption that \(A\) or \(B\) is integral.
For completeness choose sufficiently accurate integral approximations
to the selected roots and combine them by CRT. One rational variable
\(s\), chosen to be an integer in the completeness proof, works at
all the finitely many contact primes. At repeated primes choose the
same place first.

\subsection{The certificate}
\label{sec:batched-multiplication}
For nonzero \(Q_i=n_iP\), \(1\le i\le3\), write
\(x(Q_i)=x_{i0}+x_{i1}\sqrt2\). Choose rational \(a_i\)
with \(a_1a_2=a_3\), and \(d_0\in\Q^\times\). Represent
the seven scalars
\begin{equation}\label{eq:shared-coordinates}
                   d_0,\quad d_0x_{i0},\quad d_0x_{i1}
                   \quad(1\le i\le3)
\end{equation}
as \(A=(u_1+u_2+u_3)/v\). Use \(u_1,u_2,v\) as witnesses
and substitute \(u_3=Av-u_1-u_2\). Require all thirteen raw
contacts \eqref{eq:thirteen} to be nonzero, and in a signed
multiplier branch write each as \(\varepsilon kr\), where \(k\)
is from \eqref{eq:cyclic-multiplier}. There are 91 contact occurrences.

Choose one common \(J\in G\setminus\{O\}\) and one common root
approximation \(s\). At each contact choose private nonzero
\(H,X_1,X_2,X_3\in G\). Put \(Q_0=P\), \(X_0=H\),
\(Y_i=X_i-[4]Q_i\), and \(f=1/x\). Require the following
field inequalities at the place selected by the rational tests:
\begin{enumerate}[label=(\roman*),leftmargin=*,itemsep=2pt]
\item order at least one for \(f(H),f(X_1),f(X_2),f(X_3),
f(J^\sigma),f(Y_0^\sigma),\ldots,f(Y_3^\sigma)\);
\item order at least 48 for
\(\ell_{48}(z(X_i))/\ell_{48}(z(H))-a_i\), \(1\le i\le3\);
\item order at least 48 for
\(z(Y_i^\sigma)/z(J^\sigma)\), \(0\le i\le3\).
\end{enumerate}
These local inequalities are imposed by rational equations, using
\eqref{eq:batch-field-tests}
and Lemma~\ref{lem:batch} once per contact, with \(M=64\): there
are 33 rational tests, nineteen of order one and fourteen of order 48.
Every denominator is inverted. Denote the resulting existential
relation by \(\mathcal M_b\).

\begin{proposition}\label{prop:batched-error}
For nonzero integer indices, true multiplication satisfies
\(\mathcal M_b\). If \(\mathcal M_b(Q_1,Q_2,Q_3)\) holds but \(n_1n_2\ne n_3\), then
\(\max_i|n_i|\le L_*=18H_*+12\).
\end{proposition}
\begin{proof}
For soundness write \(H=hP\), \(J=jP\), \(X_i=t_iP\),
and suppose a represented slope has an odd normalized contact at
\(q\notin S_*\). Some raw contact then has odd valuation:
if all raw valuations were even, so would be their minima and
normalized depths. The multiplier is a unit at \(q\).
Lemma~\ref{lem:batch} supplies the field inequalities above.
The alternate comparisons and \eqref{eq:formal-ratio} give
\[
 v_q((h-4)/j)\ge48,\qquad v_q((t_i-4n_i)/j)\ge48.
\]
Since \(j\) is a nonzero integer, \(h\equiv4\) and
\(t_i\equiv4n_i\pmod{q^{48}}\). In particular \(h\) is a
unit. The primary comparisons and \eqref{eq:formal-log} now give
\(a_i\equiv t_i/h\equiv n_i\pmod{q^{48}}\).
Thus \(q^{48}\mid\delta:=n_1n_2-n_3\).

If \(\delta\ne0\), put \(B=\max_i|n_i|\ge1\).
For each represented slope, \(M_S(u:v)^{48}\mid\delta\),
while \(|\delta|\le2B^2\). The height bound \eqref{eq:height}
therefore gives \(h(u:v)\le2H_*B\). Each scalar in
\eqref{eq:shared-coordinates}, being a sum of three slopes, has
height at most \((6H_*+2)B\). Division by the common \(d_0\)
and recombination in \(\Q(\sqrt2)\) give
\[
 h(x(Q_i))\le3(6H_*+2)B+2.
\]
At a maximizing index, \eqref{eq:canonical} yields
\(B^2-4\le(18H_*+6)B+2\), hence \(B\le18H_*+12\).

For completeness, take \(a_i=n_i\) and an integral \(d_0\)
clearing all abscissa denominators. Apply \textbf{(I4)} to the seven
integers in \eqref{eq:shared-coordinates}, choosing a consistent
place at each distinct contact prime. For an occurrence at \(r\),
let \(h_r,j_r\) be its local witnesses. Applying the alternate
test with index 1 gives
\(h_r\equiv4\pmod{r^{48+v_r(j_r)}}\).
Choose a positive common multiple \(j\) of all nonzero \(j_r\).
The congruences
\[
 t_r\equiv1\pmod{|j_r|},\qquad
 h_rt_r\equiv4\pmod{r^{48+v_r(j)}}
\]
are compatible: their moduli have gcd \(r^{v_r(j_r)}\), and
\(4h_r^{-1}\equiv1\) modulo this gcd. The generalized CRT
gives an infinite progression, so choose a solution with
\(h_rt_r\notin\{0,4,-4\}\). Set
\[
 J=jP,\qquad H=h_rt_rP,\qquad X_i=h_rt_rn_iP.
\]
The primary points remain formal. The alternate differences are
formal because
\(h_rt_r-4=(h_r-4)+h_r(t_r-1)\), and the second term is a
multiple of \(j_r\). The strengthened congruence supplies order
48 against the common \(j\). Choose a sufficiently accurate integral
\(s\) by CRT; Lemma~\ref{lem:batch} supplies all parity witnesses.

These choices also justify the nonzero charts used in the circuit.
Equality \(x(X_i)=x([4]Q_i)\) would give
\(h_rt_rn_i=\pm4n_i\), hence \(h_rt_r=\pm4\), a contradiction;
this includes \(i=0\). All differences are nonidentity points.
A nonidentity point of the torsion-free group \(G\) has neither
zero abscissa nor zero ordinate. Finally
\(\ell_{48}(z(H))\ne0\), since at its primary formal place its
valuation equals that of the nonzero parameter. Thus every required
inverse exists. The same finite CRT argument allows \(J,s\) to be
shared across all three copies of \(\mathcal M_b\) used below.
\end{proof}

\subsection{An anchor beyond the finite error}\label{sec:anchor}
Let \(T_{2,3}(\Q)=\Z_{(2)}\cap\Z_{(3)}\), with the explicit
seven-witness formula \eqref{eq:T23}. Define the chain
\[
 c_0=3,\quad c_{i+1}=c_i^2\ (0\le i<38348),
 \qquad C=c_{38348}=3^{2^{38348}}.
\]
For \(R\in G\setminus\{O\}\), with
\(x(R)=x_0+x_1\sqrt2\), define \(\operatorname{Anchor}(R)\)
by this chain and
\begin{equation}\label{eq:anchor}
\begin{gathered}
 x_0f_0+2x_1f_1=1,\qquad x_0f_1+x_1f_0=0,\\
 Ct_0=f_0,\qquad Ct_1=f_1,\qquad
 T_{2,3}(t_0)\ \wedge\ T_{2,3}(t_1).
\end{gathered}
\end{equation}
The constant \(C\) is represented by the chain, not by its expansion.

\begin{lemma}\label{lem:anchor-exists}
Anchor points exist, and every anchor \(R=sP\) satisfies
\(|s|>L_*\). An anchor with positive index exists.
\end{lemma}
\begin{proof}
The equations give \(v_3(1/x(R))\ge2^{38348}\).
The local calculation in Appendix~\ref{sec:anchor-certificate}
gives \(v_3(1/x(sP))=2+2v_3(s)\), hence
\(\lvert s\rvert\ge3^{2^{38347}-1}>L_*\).
That appendix also verifies the last inequality and constructs anchors
by taking multiples deep in the formal groups above 2 and 3.
Negation preserves their abscissas, so the index can be positive.
\end{proof}

\begin{proposition}[Exact multiplication using one anchor]
\label{prop:anchor-multiplication}
For nonzero \(A=aP,B=bP,C'=cP\), the relation \(ab=c\)
is equivalent to
\begin{equation}\label{eq:anchor-multiplication}
\begin{split}
 \exists R,U,V\in G\setminus\{O\}:\quad
 &\operatorname{Anchor}(R)\ \wedge\ \mathcal M_b(R,A,U)\\
 &\wedge\ \mathcal M_b(R,C',V)\ \wedge\ \mathcal M_b(U,B,V).
\end{split}
\end{equation}
\end{proposition}
\begin{proof}
Write \(R=sP\). The first two approximate multiplications are
exact, since \(|s|>L_*\). Thus \(U=saP\) and \(V=scP\).
Since \(a\ne0\), also \(|sa|>L_*\), so the last one is
exact and gives \((sa)b=sc\), or \(ab=c\). Conversely choose
any anchor and these true products. For arbitrary indices define
\[
\begin{split}
\Mult(A,B,C):\quad&G(A)\wedge G(B)\wedge G(C)\ \wedge\\
 &\bigl[((A=O\vee B=O)\wedge C=O)\\
 &\qquad\vee(A,B,C\ne O\wedge\text{\eqref{eq:anchor-multiplication}})\bigr].
\end{split}
\]
This includes exactly the zero products. There is no use of index multiplication in the anchor
equations themselves.
\end{proof}

\subsection{Point codes and the multiplication polynomial}\label{sec:full-multiplication}

The completeness proof of Proposition~\ref{prop:batched-error}
justifies nonzero affine private points and the ordinary subtraction
chart. Every private nonzero point of \(G\) is generated as \(2T\) from an
affine curve point \(T\) with \(y(T)\ne0\), using the ordinary
doubling chart. Conversely every nonzero point of \(2E(F)\) has such
a preimage. The torsion of \(E(F)\) has exponent two, by
Appendix~\ref{sec:numerical-elliptic}. Consequently the required inverse abscissas and parameters are
well-defined on these nonzero points.

For the three external inputs use codes
\((\tau,x_0,x_1,y_0,y_1)\), with \(\tau\in\{0,1\}\).
Impose \((1-\tau)x_i=(1-\tau)y_i=0\), and let
\(\bar A\) be the affine point whose coordinates are
\(\tau_A(x_A,y_A)+(1-\tau_A)(x_P,y_P)\), and similarly for
\(\bar B,\bar C\). Attest each barred point as a double. This
allows precisely the code of \(O\) and the affine nonzero points
of \(G\). Set \(\epsilon=\tau_A\tau_B\), require \(\tau_C=\epsilon\), and apply the nonzero multiplication circuit to the three affine points
with coordinates
\[
 \epsilon\bar A+(1-\epsilon)P,\quad
 \epsilon\bar B+(1-\epsilon)P,\quad
 \epsilon\bar C+(1-\epsilon)P,
\]
where these expressions denote coordinate interpolation by a Boolean
variable. If \(\epsilon=0\), the inner triple is \((P,P,P)\)
and the output code is \(O\). If \(\epsilon=1\), the inner
triple is the actual input triple. This handles all zero cases with one nonzero circuit.


The field operations are
\((a,b)(c,e)=(ac+2be,ae+bc)\), conjugation \((a,b)\mapsto(a,-b)\),
and inversion \((a,b)^{-1}=(a,-b)/(a^2-2b^2)\).
All divisions use inverse witnesses. For points \(U,V,W\),
the addition law consists of the identity and opposite-point cases,
or, when all are affine,
\[
 x_W=m^2-x_U-x_V,\qquad y_W=m(x_U-x_W)-y_U,
\]
with \(m(x_V-x_U)=y_V-y_U\) if \(x_U\ne x_V\), and
\(2y_Um=3x_U^2-d^2\) if \(U=V\), \(y_U\ne0\).
Splitting in the basis \(1,\sqrt2\) gives rational equations.
Every gate has a quadratic residual; summing their squares preserves
conjunction over \(\Q\).

The program \texttt{compile\_multiplication.py} supplies every field
operation, point operation, formal-logarithm evaluation, raw contact,
multiplier choice, and rational test in the three copies of
\(\mathcal M_b\).  One \(J\) and one root approximation are shared
across all 273 contacts. The common-multiple and CRT proof in
Proposition~\ref{prop:batched-error} applies to this finite collection;
the primary contact witnesses remain private.

Clear denominators in each quadratic residual and sum their squares.
The resulting integer polynomial \(P_{\rm mult}\) has degree four,
15 free coordinates, 511818 rational witnesses and 509019 residuals.
By Proposition~\ref{prop:anchor-multiplication},
\[
 \exists\mathbf w\in\Q^{511818}\quad
 P_{\rm mult}(\code(A),\code(B),\code(C),\mathbf w)=0
\]
holds exactly when there are integers \(a,b,c\) with
\(A=aP\), \(B=bP\), \(C=cP\) and \(ab=c\).
Thus invalid point codes are rejected.
The complete leaf circuit is included in the accompanying files.
Its input slots are the coordinates of \(A,B,C\), in that order;
each point is ordered as \((\tau,x_0,x_1,y_0,y_1)\).

\section{The integrality formula}\label{sec:integers}
The following construction turns multiplication on \(G\) into a
definition of the literal subset \(\Z\subset\Q\), following
\cite[Sections 5--7]{Conant}.
Set
\[
 \beta(t)=\frac{t}{(1+t^2)^2},\quad
 c_n=\Norm_{F/\Q}(z(nP)),\quad b_n=\beta(c_n),\quad
 R(n)=\prod_{p\in\Supp(b_n)}p\qquad(n\ge1).
\]
The support includes primes dividing the reduced numerator or denominator.
All \(b_n\) are nonzero.

\subsection{A support bound and a bounded divisibility relation}
\begin{lemma}\label{lem:support-majorant}
For every positive integer \(n\),
\[
                         R(n)<2^{1024n^2}.
\]
\end{lemma}
\begin{proof}
On the scaled curve \(E_0\) of Appendix~\ref{sec:numerical-elliptic},
\(h(x(P_0))=\log41\). Lemma~\ref{lem:duplication-height} implies
\[
 \mathfrak h(P_0)\le\log41+\tfrac13\log72<\tfrac{11}{2},
 \qquad
 h(x(nP_0))<\tfrac{11}{2}n^2+2.
\]
The coordinate change has height less than 2. Hence
\(h(x(nP))<\tfrac{11}{2}n^2+4<10n^2\).
The element \(d\) is integral, with one positive conjugate below
1 and the other below 150, so \(h(d)<\tfrac12\log150<3\).
Writing \((x,y)=nP\), the curve equation and elementary height
inequalities give
\[
 h(y)\le2h(x)+h(d)+\tfrac12\log2<2h(x)+\tfrac72,
 \qquad h(z(nP))\le h(x)+h(y)<34n^2.
\]
Height is invariant under \(\sigma\), so
\(h(c_n)\le2h(z(nP))<68n^2\). If \(c_n=a/b\) in lowest
terms, \(b>0\), then
\[
 \beta(c_n)=\frac{ab^3}{(a^2+b^2)^2},\qquad
 h(b_n)\le4h(c_n)+\log4<274n^2.
\]
Let \(u/v\) be the reduced expression for \(b_n\), with
\(v>0\). Since \(u\ne0\),
\[
 R(n)\le |u|v\le\exp(2h(b_n))
       <\exp(548n^2)<3^{548n^2}<2^{1024n^2}.
\]
The numerical comparisons follow from the elementary exponential
bounds recorded in Appendix~\ref{sec:saturation}.
\end{proof}

For positive \(n,s\), put
\begin{equation}\label{eq:majorant-bound}
 B=(s^2+1)4^{12288n^2},\qquad Q=4^B,\qquad T=4Q^2.
\end{equation}
For positive integers \(T,\mathcal F,\lambda\), define
\begin{equation}\label{eq:divisibility-replacement}
 \mathcal D(T,\mathcal F,\lambda):\quad
 \forall k\ (1\le k\le T\ \Longrightarrow\
 k\mid \mathcal F\ \wedge\ [k=1\ \vee\ k\text{ composite}\ \vee\
                  (k\mid\lambda\ \wedge\ k^2\nmid\lambda)]).
\end{equation}
This bounded natural-number specification will be replaced by explicit
existential equations in Section~\ref{sec:completed-definition}.

\begin{lemma}\label{lem:divisibility-replacement}
Condition \eqref{eq:divisibility-replacement} is equivalent to
\[
 \operatorname{lcm}(1,\ldots,T)\mid \mathcal F,\qquad
 v_p(\lambda)=1\quad\text{for every prime }p\le T.
\]
Such positive \(\mathcal F,\lambda\) exist for every positive \(T\).
\end{lemma}
\begin{proof}
The divisibility conditions on \(\mathcal F\) give the least common multiple.
At a prime \(k=p\), the first two alternatives are false, leaving
exactly \(p\mid\lambda\), \(p^2\nmid\lambda\); at 1 or a
composite \(k\) another alternative applies. Existence follows by
taking \(\mathcal F=T!\), \(\lambda=\prod_{p\le T}p\).
\end{proof}

Let \(\Gamma^{\rm div}(n,s,M,N)\) assert \(n,s\ge1\) and
the existence of positive \(\mathcal F,\lambda\) satisfying
\eqref{eq:majorant-bound}, \(\mathcal D(T,\mathcal F,\lambda)\),
\(M=\lambda \mathcal F\), \(N=\lambda M\). Write
\(\widehat\Gamma^{\rm div}(A,S,U,V)\) for its interpretation
in \(G\), meaning \((A,S,U,V)=(nP,sP,MP,NP)\).
Section~\ref{sec:completed-definition} constructs this existential
interpretation using the multiplication polynomial, addition and four
squares on indices. Factorials and prime products are used only as
existence witnesses in the lemma, not as operations in that construction.

\subsection{Definition and proof}
For nonidentity point codes put
\[
 b(A)=\beta(\Norm_{F/\Q}(z(A))),\qquad
 t(U,V)=\tfrac12\Tr_{F/\Q}(z(V)/z(U)).
\]
Using the local formulas of Appendix~\ref{sec:local-rings}, define
\begin{equation}\label{eq:integer-formula}
\begin{split}
 \mathcal Z(x):\quad&x(x-1)(x-2401)=0\quad\vee\\
 &\Bigl[\bigwedge_{p\in S_*}\Int_p(x)\ \wedge\
 \exists A,S,U,V\ \bigl(\operatorname{Anchor}(S)\ \wedge\
 \widehat\Gamma^{\rm div}(A,S,U,V)\\
 &\qquad\wedge\Phi(x)\ \wedge\
 \bigwedge_{j=1}^4\bigl(L_j(x,1)\ne0\ \wedge\
       \Phi(b(A)/L_j(x,1))\bigr)\\
 &\qquad\wedge\Phi(\beta(t(U,V))x^2)\bigr)\Bigr].
\end{split}
\end{equation}
Rational functions are expanded using equations and inverse witnesses.

\begin{proposition}\label{prop:integers}
The formula \(\mathcal Z\) defines \(\Z\) in \(\Q\).
\end{proposition}
\begin{proof}
Suppose the main disjunct holds, with \(A=nP\), \(S=sP\).
An odd contact prime \(q\notin S_*\) cannot be an odd denominator
prime of \(x\), by \(\Phi(x)\). An even denominator valuation
gives even depth at infinity and zero finite contact depths, so it
cannot occur at this \(q\) either. Thus \(x\) is integral at
\(q\), and some \(L_j(x,1)\) has positive odd valuation.
The test \(\Phi(b_n/L_j(x,1))\) forces \(q\in\Supp(b_n)\).
The contact-height bound and Lemma~\ref{lem:support-majorant} imply
\[
 h(x)\le H_*R(n)^{24}
       <H_*2^{24576n^2}
       <(s^2+1)4^{12288n^2}=B,
\]
because the positive anchor index satisfies \(s>L_*>H_*\).

For a prime \(p\le Q\) outside \(S_*\), every reduction group
above \(p\) has order at most \(4Q^2=T\), by Hasse's bound and
\([F:\Q]=2\). Its order divides \(\mathcal F\), so \(\mathcal F P\) is formal
at every such place. Since \(v_p(\lambda)=1\), the parameter of
\(MP=\lambda\mathcal F P\) has valuation at least two. The integral formal
multiplication law gives
\[
 \frac{z(NP)}{z(MP)}\equiv\lambda\pmod{p^2\mathcal O_v},
 \qquad v_p(t(U,V))=1.
\]
For the second assertion take trace and divide by two. A denominator
prime \(p\) of \(x\) would lie outside \(S_*\) by the local
tests, and satisfy \(p\le e^{h(x)}<e^B<4^B=Q\). Then
\[
 v_p(\beta(t(U,V))x^2)=1+2v_p(x)<0
\]
is odd, contradicting the last parity test. Hence \(x\in\Z\).

Conversely let \(x\in\Z\setminus\{0,1,2401\}\). Choose a
positive multiple \(nP\) sufficiently close to \(O\) at all
places above the primes dividing \(\prod_jL_j(x,1)\), so that
\(v_p(c_n)\ge\max_jv_p(L_j(x,1))>0\) there. Such multiples
exist: small formal neighborhoods are finite-index subgroups of the
compact local elliptic groups, and their preimages in \(\Z P\)
have nonzero intersection. The same bound holds for \(b_n\).
Elsewhere every negative valuation of \(b_n/L_j(x,1)\) is even,
since for coprime integers \(a,b\),
\[
                 \beta(a/b)=\frac{ab^3}{(a^2+b^2)^2}.
\]
Thus all four ratios have square reduced denominator. The same is
true of \(\beta(t)x^2\) for every rational \(t\).
Choose a positive anchor index \(s\); this is possible because the
anchor depends only on the abscissa and is invariant under negation.
The lemma supplies \(\mathcal F,\lambda\), hence \(U,V\).
Proposition~\ref{prop:explicit-phi} supplies every parity witness;
the local tests hold for integers. The exceptional disjunct adds
precisely the integral contact values \(0,1,2401\).
\end{proof}

\section{The explicit polynomial}\label{sec:completed-definition}
We now eliminate the bounded quantifier and specify the resulting
polynomial by fixed finite blocks with fresh variables.

\subsection{A small elementary matrix for the bounded condition}

Put \(k=i+1\), where \(0\le i<T\). Quantify nonnegative
\(d,b_0,b_1,b_2,a,b,u,v,r,w\), and impose
\begin{align}
 \mathcal F&=kd,& b_0+b_1+b_2&=1,& b_0i&=0,\nonumber\\
 b_1(a+2)(b+2)&=b_1k,&
 b_2\lambda&=b_2ku,\label{eq:elementary-divisibility}\\
 b_2u&=b_2(kv+r+1),&
 b_2k&=b_2(r+2+w).\nonumber
\end{align}
The selectors are automatically zero or one because they are natural
numbers with sum one. The first branch is \(k=1\); the second
provides a proper factorization; the last expresses
\(\lambda=ku\) with nonzero remainder \(r+1\) of \(u\) modulo
\(k\). Its inequality is exactly \(1\le r+1<k\).
Thus this is equivalent to the matrix of
\eqref{eq:divisibility-replacement}. Introducing ordinary arithmetic
gates gives 31 natural witnesses and 28 equations, each of the form
\(z=x+y\) or \(z=xy\), allowing fixed constants and repeated
outputs. The file \texttt{natural/body.json} specifies every equation.

\subsection{The power-of-four component}\label{sec:executed-dprm}
The constructive DPRM development \cite{DPRMCoq}, in the pinned
Coq source \cite{DPRMLibrary}, supplies an extracted
finite list of natural-number equations for \(Q=4^B\). The wrapper
\texttt{power4\_constraints\_correct} proves
\begin{equation}\label{eq:power4-checked}
 Q=4^B\quad\Longleftrightarrow\quad
 \exists\mathbf w\in\N^{7447}\ \bigwedge_{e=1}^{5780}r_e(Q,B,\mathbf w)=0.
\end{equation}
Each residual is a constant, parameter, alias, addition or multiplication
equation. Coq 8.18 accepted the wrapper without additional assumptions.
Eliminating aliases and operations with zero or one gives an equivalent
list of 5346 natural witnesses and 3679 equations. The supplied
reconstruction map recovers every eliminated witness; a separate replay
substitutes it into every original equation. Two copies, linked by
\(D=4^{12288n^2}\), \(B=(s^2+1)D\) and \(Q=4^B\),
give exactly the required bounds.

\subsection{Eliminating the bounded quantifier by explicit templates}

We use the sequence encoding proved in the pinned Coq development
\cite{DPRMLibrary}, files \texttt{Matija/cipher.v},
\texttt{Dio/dio\_cipher.v}, and \texttt{Dio/dio\_bounded.v}.
For a tuple \(a_0,\ldots,a_{T-1}\) with \(a_i<2^q\) and
\(q>T+1\), its code is
\[
 \operatorname{enc}_{T,q}(a)
       =\sum_{i<T}a_i\rho^{2^{i+1}},\qquad \rho=2^{4q}.
\]
The relevant polynomial relations are \(\operatorname{Code}\),
asserting that an integer is such a code; \(\operatorname{CodeNat}\),
encoding \(0,\ldots,T-1\); \(\operatorname{Const}\), encoding a
constant tuple; and \(\operatorname{CodeMult}\), imposing componentwise
multiplication on three codes already known to be valid. Componentwise
addition is ordinary addition of codes. The cited development proves
these properties, including the absence of carries in the required masks.

Choose \(q=T+2+h\), \(h\in\N\). Encode the index by
\(\operatorname{CodeNat}\), the two parameters \(\mathcal F,\lambda\)
and constants \(0,1,2\) by \(\operatorname{Const}\), and each
of the 31 matrix witnesses by \(\operatorname{Code}\). Replace
every addition equation by the addition of its codes and every
multiplication equation by \(\operatorname{CodeMult}\).
This uses precisely one \(\operatorname{CodeNat}\), five
\(\operatorname{Const}\), 31 \(\operatorname{Code}\), and
13 \(\operatorname{CodeMult}\) instances.

For soundness, decode each tuple. The code relations imply every
matrix equation at each \(i<T\), giving the bounded condition.
For completeness, choose a witness to the matrix for each of the
finitely many \(i<T\); then choose \(q\) larger than \(T+1\)
and all their binary lengths, and encode these tuples. This proves
the equivalence without retaining a bounded quantifier in the final formula.

Coq 8.18 checks all four template constructions without additional
assumptions. After alias elimination their sizes are:
\begin{center}\small
\begin{tabular}{lrr}\toprule
Template & Natural witnesses & Elementary equations\\\midrule
\(\operatorname{Code}\)&138563&95337\\
\(\operatorname{CodeNat}\)&459906&316434\\
\(\operatorname{Const}\)&121966&83918\\
\(\operatorname{CodeMult}\)&221849&152642\\
Power of four&5346&3679\\\bottomrule
\end{tabular}
\end{center}
Reconstruction and substitution into every original equation verify
the alias elimination for each template.

\subsection{The full rational-field interpretation}

Every natural variable becomes a five-coordinate elliptic point code.
Require its index to be a sum of four squares, using four copies of
\(P_{\rm mult}\) and three elliptic additions. This gives the
polynomial block \texttt{q\_natural}, with 2048167 rational witnesses.
Each elementary integer addition or multiplication is replaced by the
corresponding rational point-code block. Every free natural input is
restricted at the outer graph block; every local natural witness is
restricted in the block that introduces it. Thus no unrestricted rational
variable is silently treated as a natural number.

The addition block validates all three elliptic point codes. Five
Boolean selectors, with sum one, choose: the first input is \(O\);
the second input is \(O\); the inputs are opposites; ordinary doubling;
or ordinary unequal-abscissa addition. All necessary inverse conditions
are enforced in the selected branch. In an inactive branch its denominator
is replaced by one, so no extraneous exclusion is imposed. The block has
281 rational witnesses and 342 quadratic residuals. It expresses the
usual elliptic group law; when its inputs lie in \(G\), so does its output.
Together with the already computed multiplication block and Lagrange's
four-square theorem, this justifies the interpretation of every natural
constraint. Fixed integer indices such as 12288 are constructed by binary
elliptic addition; their enormous point coordinates are not expanded.

The root block has the single free rational input \(x\). A Boolean
selector \(b\) imposes
\[
 (1-b)x(x-1)(x-2401)=0,\qquad x'=bx+2(1-b),
\]
using quadratic gates. All remaining conditions are the main integrality
disjunct applied to \(x'\), with
\(\widehat\Gamma^{\rm div}(A,S,U,V)\) and
\(\operatorname{Anchor}(S)\). When \(b=0\), the input is an
integral exceptional contact and \(x'=2\). When \(b=1\), the actual
input is tested. This avoids duplicating the entire construction for the
three exceptional values.

All 7499 primes at most 76176 are enumerated, and their fixed local-ring
formulas from Appendix~\ref{sec:local-rings} are emitted. For the
\(p\equiv1\pmod8\) cases the largest auxiliary prime required is
131; all primality and Legendre-symbol conditions are checked exactly.
The six parity inputs are \(x'\), the four contact ratios, and
\(\beta(t(U,V))x'^2\). The nonzero parity block suffices here:
its equation forces \(x'\ne0\), the contact denominators are
explicitly inverted, and \(b(A)\ne0\). Also \(t(U,V)\ne0\),
because the graph gives valuation one at every prime outside \(S_*\)
and below \(Q\). Such a prime exists: \(Q\) exceeds the first
prime above 76176, by its displayed exponential bound (or Bertrand's
postulate). Thus the nonzero restriction preserves completeness.

\subsection{An exact finite polynomial specification}

The file \texttt{Z\_definition.polynomial.json} and its listed
dependencies specify the polynomial as follows. Each block
\(\mathcal B\) has a fixed list of free input slots, local variables,
integer-coefficient quadratic residuals \(r_j\), and calls to other
blocks. Its polynomial is
\begin{equation}\label{eq:modular-polynomial-semantics}
 P_{\mathcal B}=\sum_j r_j^2+
       \sum_{\text{calls }c}\ \sum_{i=0}^{m_c-1}
                  P_{\mathcal B_c}(\alpha_{c,i};\mathbf w_{c,i}).
\end{equation}
Every \(m_c\) is a displayed fixed positive integer, every input
substitution \(\alpha_{c,i}\) is an explicitly listed variable slot
or an explicitly listed affine progression of slots, and
\(\mathbf w_{c,i}\) is a fresh block of witnesses. Calls form an
acyclic finite dependency graph. Consequently
\eqref{eq:modular-polynomial-semantics} is an exact specification of a
single polynomial in finitely many indeterminates. It contains no
arithmetic predicate, unspecified constant, or unexecuted DPRM conversion.
The repetitions are fixed finite sums in the polynomial's notation,
not quantifiers in the defining formula.

The residuals of the multiplication leaf are stored completely in the
included line-based circuit file; every other leaf is directly listed
in its block file. All rational coefficients were cleared in each residual.
All blocks are sums of squares; therefore their sum is zero over \(\Q\)
if and only if every constituent constraint is zero. Input substitutions
only rename variables. Degree remains at most four, and a private squaring
residual shows that it is exactly four.

The exact output counts, including fresh witnesses in every repeated
block, are
\begin{center}
\begin{tabular}{lr}\toprule
Quantity & Value\\\midrule
Free rational input &1\\
Existential rational witnesses &\(17\,030\,326\,787\,289\)\\
Primitive squared residuals &\(16\,938\,632\,062\,448\)\\
Degree &4\\
Copies of the multiplication polynomial &\(33\,256\,658\)\\\bottomrule
\end{tabular}
\end{center}
Denote this particular polynomial by \(P_{\rm final}\).

\begin{proof}[Proof of Theorem~\ref{thm:main}]
The decoded bounded certificate is equivalent to
\eqref{eq:divisibility-replacement}. Its interpreted point-code equations
are equivalent to the natural equations, by the multiplication theorem,
the addition block and four squares. Proposition~\ref{prop:integers}
therefore proves the main integrality branch. The selector restores
precisely the integral exceptional contacts. Finally the sum-of-squares
identity \eqref{eq:modular-polynomial-semantics} combines exactly those
equations and gives the stated degree and witness count.
\end{proof}

\paragraph{Polynomial files and verification scope.}
The accompanying \texttt{completed\_polynomial.zip} contains
\texttt{Z\_definition.polynomial.json}, all block and leaf files,
source code, template proofs and execution records. The dependency
fingerprint (SHA-256 of the sorted list of block and leaf hashes) is
\begin{center}\small\ttfamily
7cff5a2443678289d21a468d7a58e683\\
227a7314ee31f520157c359aa6493c4f
\end{center}
with the two lines concatenated. Coq checks the natural-number
constructions; extraction, serialization, optimization and the
rational-field translation have separate exact checks. These do not
constitute a formal verification of the complete proof.

\appendix

\section{Numerical elliptic data}\label{sec:numerical-elliptic}

\subsection{Torsion and the displayed point}

Counting points on the two reductions at 17, and on a reduction at 23,
gives the following exact table. The middle column specifies the image
of \(\sqrt2\) in the residue field.
\[
\begin{array}{c|r|r|r}
p&\sqrt2&d&\#\widetilde E(\F_p)\\\hline
17&6&12&20\\
17&11&2&16\\
23&5&17&24\\
23&18&18&24
\end{array}
\]
The counts are the elementary sums
\(p+1+\sum_{x\in\F_p}\chi_p(x^3-d^2x)\), with \(\chi_p(0)=0\).
The prime-to-17 torsion injects into both groups at 17, so its order
divides 4. Its 17-primary part injects at 23 and is trivial. The
four visible points \(O,(0,0),(d,0),(-d,0)\) exhaust the torsion.
Consequently
\[
 E(F)_{\rm tors}=E[2](F)\cong(\Z/2\Z)^2,\qquad m=2.
\]
By \cite[Proposition 2.1, proved in Section 5]{H10}, this exact
curve has rank one. Hence \(2E(F)\) is infinite cyclic.

A particularly simple infinite-order point and its double are
\begin{align*}
R&=(-7+5\sqrt2,-280+198\sqrt2),\\
P=2R
 &=\left(\frac{1641}{8}-145\sqrt2,
             -\frac{15465}{4}+\frac{87483}{32}\sqrt2\right).
\end{align*}
The coordinates satisfy the curve equation exactly, and \(R\) is
nontorsion because its second coordinate is nonzero. Under
\((x,y)\mapsto(v^2x,v^3y)\), where \(v=3+2\sqrt2\), the curve
becomes
\[
 y^2=x^3-(3-\sqrt2)^2x,
 \qquad R\mapsto(1+\sqrt2,2\sqrt2).
\]
Write this model as \(E_0\), put \(d_0=3-\sqrt2\), and put
\(D_0=d_0^2=11-6\sqrt2\). Denote the images of \(R,P\) by
\(R_0,P_0\). We verify primitivity and the numerical height bounds below.

\subsection{A numerical duplication-height estimate}
\label{sec:duplication-height}

Use the absolute logarithmic Weil height \(h\). For \(T\in E_0(F)\)
define
\[
 \mathfrak h(T)=\lim_{j\to\infty}4^{-j}h(x(2^jT)),
 \qquad h(x(O))=0.
\]
This is twice the usual N\'eron--Tate height; in particular
\(\mathfrak h(nT)=n^2\mathfrak h(T)\), and adding torsion does
not change it.

\begin{lemma}\label{lem:duplication-height}
For every \(T\in E_0(F)\),
\[
 4h(x(T))-\log392\le h(x(2T))\le4h(x(T))+\log72.
\]
Consequently
\begin{equation}\label{eq:explicit-height-comparison}
 h(x(T))-\tfrac13\log392\le\mathfrak h(T)
       \le h(x(T))+\tfrac13\log72.
\end{equation}
\end{lemma}
\begin{proof}
The duplication morphism on the projective \(x\)-line is
\[
 F_1(X,Z)=(X^2+D_0Z^2)^2,\qquad
 F_2(X,Z)=4XZ(X^2-D_0Z^2).
\]
At each place normalize \(\max(|X|,|Z|)=1\) and compare
\(\max(|F_1|,|F_2|)\) with 1. Both real values of \(D_0\)
exceed 1. For real \(t\),
\((t^2+1)^2\ge4|t(t^2-1)|\), as follows by expanding
\((t^2-2t-1)^2\) and \((t^2+2t-1)^2\).
Substitution \(t=X/(\sqrt{D_0}Z)\) therefore gives
\[
 1\le F_1(X,Z)=\max(|F_1(X,Z)|,|F_2(X,Z)|)
       \le(1+D_0)^2.
\]
The case \(Z=0\) is immediate. The sum of the two archimedean
contributions, each weighted by \(1/2\), is between 0 and
\(\log((1+D_0)(1+D_0^\sigma))=\log72\).

At a finite prime use integral coordinates with
\(\min(v(X),v(Z))=0\). Outside 2 and the prime
\(\mathfrak d=(3-\sqrt2)\) of norm 7, the two homogeneous
polynomials have no common zero over the residue field. Thus the
local loss is zero. At \(\mathfrak d\), one has \(v(D_0)=2\).
If \(X\) is a unit, \(F_1\) is a unit. Otherwise \(Z\) is a
unit. If \(v(X)>1\), then \(v(X^2+D_0Z^2)=2\); if
\(v(X)=1\), the same equality follows because \(-1\) is not
a square in \(\F_7\). Hence \(v(F_1)\le4\), giving a height
loss at most \(\tfrac42\log7\).

At \(\pi=\sqrt2\), normalize \(v_\pi(\pi)=1\), so
\(v_\pi(2)=2\). The element \(D_0\) is a unit. If exactly
one of \(X,Z\) is a unit, \(F_1\) is a unit. If both are
units, write \(X/(d_0Z)=a+b\pi\), with \(a\in\Z_2^\times\)
and \(b\in\Z_2\). Expansion shows
\[
 v_\pi((a+b\pi)^2+1)=
 \begin{cases}2,&b\in2\Z_2,\\3,&b\notin2\Z_2.\end{cases}
\]
Indeed, the rational term is \(a^2+1+2b^2\), and the other
term is \(2ab\pi\). Thus \(v_\pi(F_1)\le6\), giving a
height loss at most \(\tfrac62\log2\).
The total finite loss is at most
\(2\log7+3\log2=\log392\), proving the first statement.
The second follows by telescoping and summing
\(\sum_{j\ge0}4^{-j-1}=1/3\).
\end{proof}

\subsection{A finite primitivity certificate}
\label{sec:saturation}

\begin{proposition}\label{prop:generator}
The point \(R\) generates \(E(F)\) modulo torsion, and the displayed
point \(P=2R\) generates \(2E(F)\).
\end{proposition}
\begin{proof}
The free part has rank one. The two-isogeny homomorphism
\(\alpha:E_0(F)\to F^\times/F^{\times2}\), given by
\(x\) off \(O,(0,0)\), by 1 at \(O\), and by \(-D_0\)
at \((0,0)\), kills \(2E_0(F)\); see \cite[Section 5]{H10}.
All torsion images have equal signs at the two real embeddings,
whereas \(x(R_0)=1+\sqrt2\) has opposite signs. Thus
\(R_0\notin2E_0(F)+E_0(F)_{\rm tors}\), so its index in the
free part is odd.

The upper bound in \eqref{eq:explicit-height-comparison} gives
\[
 \mathfrak h(R_0)
 \le\tfrac12\log(1+\sqrt2)+\tfrac13\log72<2.
\]
If \(R_0=rT+\tau\) for a generator \(T\) of the free part
and a torsion point \(\tau\), with \(|r|>1\), then \(|r|\ge3\).
It would follow that
\[
 \mathfrak h(T)<\tfrac29,\qquad
 h(x(T))<\tfrac29+\tfrac13\log392<\tfrac{20}9<\tfrac94.
\]
We exhaust a finite set containing all possible abscissas of such a
point \(T\), as follows.

If \(x\in\Q\) and \(h(x)<9/4\), write it in reduced form
\(u/v\), where \(|u|\le9\) and \(1\le v\le9\), since
\(e^{9/4}<10\). If \(x\in F\setminus\Q\), its primitive
minimal polynomial \(aX^2+bX+c\), with \(a>0\), has Mahler
measure \(M=e^{2h(x)}<e^{9/2}<181/2\). The coefficient bounds
\(a,|c|\le M\), \(|b|\le2M\) give
\begin{equation}\label{eq:search-box}
 1\le a\le90,\qquad |b|\le180,\qquad |c|\le90.
\end{equation}
Such an irreducible polynomial splits in \(F\) precisely when
\(b^2-4ac=2k^2>0\) for an integer \(k\). There are exactly
8356 primitive triples in this box with that property.

For both roots of each of these quadratics, and for every rational
candidate, test whether \(x^3-D_0x\) is a square in \(F\).
This exact enumeration yields the three torsion abscissas
\(-3+\sqrt2,0,3-\sqrt2\), and precisely these four others:
\[
\begin{array}{c|c}
x&\text{one value of }y\\\hline
1+\sqrt2&2\sqrt2\\
2+4\sqrt2&8+9\sqrt2\\
5/2-2\sqrt2&3-11\sqrt2/4\\
23-17\sqrt2&160-114\sqrt2
\end{array}
\]
The accompanying standard-library Python program reproduces the
entire enumeration. Its square test is exact: for \(z=A+B\sqrt2\)
with \(B\ne0\), test whether \(A^2-2B^2\) is a rational square
\(n^2\), and whether either \((A+n)/2\) or \((A-n)/2\) is a
nonzero rational square \(r^2\). A square root is then
\(r+(B/(2r))\sqrt2\). If \(B=0\), test whether \(A\) or
\(A/2\) is a rational square. These tests are necessary and
sufficient by squaring \(r+t\sqrt2\).

All four nontorsion abscissas have the same doubled abscissa,
independently of the sign of \(y\):
\[
 x(2T)=\frac{57}{8}-\frac72\sqrt2.
\]
Its primitive minimal polynomial is \(64X^2-912X+1681\).
Both roots exceed 1, so \(h(x(2T))=\tfrac12\log1681=\log41\).
The lower bound in \eqref{eq:explicit-height-comparison} now gives
\[
 \mathfrak h(T)\ge
 \tfrac14\bigl(\log41-\tfrac13\log392\bigr)>\tfrac14,
\]
contradicting \(\mathfrak h(T)<2/9\). This proves primitivity.

All elementary exponential inequalities used to select the search
box can also be checked rationally: the program encloses \(e\)
between \(\sum_{j=0}^7 1/j!\) and that sum plus
\(9/(8\cdot8!)\), then compares integer powers. No floating-point
height estimate is used in the certificate.
\end{proof}

\begin{corollary}\label{cor:height-constants}
For the displayed point \(P\) on the original model,
\[
                 h(x(nP))\ge n^2-4\qquad(n\ne0).
\]

\end{corollary}
\begin{proof}
The preceding computation applied to \(P_0=2R_0\) gives
\(\mathfrak h(P_0)\ge\log41-\tfrac13\log392>1\).
The upper comparison in \eqref{eq:explicit-height-comparison}
therefore implies
\(h(x(nP_0))>n^2-\tfrac13\log72>n^2-3/2\).
The change of coordinates has factor
\(v^2=(1+\sqrt2)^4\), whose height is
\(2\log(1+\sqrt2)<2\). Using
\(h(ab)\le h(a)+h(b)\) gives
\(h(x(nP))>n^2-7/2\), which implies the asserted weaker bound.
\end{proof}

\subsection{Equations for the elliptic domain}

The domain \(G=2E(F)\) itself is now explicit. For an affine point
\(A=(X,Y)\), put \(D=d^2\) and quantify \(u,v,w\in F\):
\begin{align*}
 v^2&=u^3-Du,&vw&=1,\\
 4v^2X&=(u^2+D)^2,&
 8v^3Y&=u^6-5Du^4-5D^2u^2+D^3.
\end{align*}
Adjoin the disjunct \(A=O\). These are precisely the doubling
equations, and splitting each equality in the basis \(1,\sqrt2\)
gives equations over \(\Q\).

\subsection{The truncated logarithm}

With \(z=-x/y\) and \(w=-1/y\), write \(w=z^3f\). Then
\(f=1-d^2z^4f^2\), and the solution with constant term 1 is
\[
 f=\frac{2}{1+\sqrt{1+4d^2z^4}}.
\]
Since \(x=z/w\), differentiation yields
\[
 \frac{dx}{2y}
 =\frac12\left(\frac{zw'}w-1\right)dz
 =\frac{dz}{\sqrt{1+4d^2z^4}}.
\]
The normalized formal logarithm is therefore
\begin{equation}\label{eq:closed-log}
 \ell(z)=\sum_{j\ge0}
       \frac{(-1)^j\binom{2j}{j}}{4j+1}\,d^{2j}z^{4j+1}.
\end{equation}
Thus \(\ell_{48}\) is the sum in \eqref{eq:closed-log} for
\(0\le j\le11\). Multiplication by
\(3\cdot25\cdot13\cdot17\cdot29\cdot37\cdot41\) clears its
rational denominators without changing a quotient of two evaluations.
The conjugate polynomial replaces \(d\) by \(75+53\sqrt2\).

\subsection{The anchor certificate}\label{sec:anchor-certificate}

The prime 3 is inert in \(F\), and
\(\mathcal O_{F_v}=\Z_3[\sqrt2]\). With \(v(3)=1\),
\(v(a+b\sqrt2)=\min(v_3(a),v_3(b))\).
The displayed model of \(E\) has good reduction at 3 because
\(d\) and its discriminant are units there. Direct exact arithmetic gives
\[
 P\bmod3=(2\sqrt2,0),\qquad
 z(2P)=\frac{-6351748793760-5612900635092\sqrt2}{3287811625151}.
\]
The reduction of \(P\) has order two; the denominator in the last
expression is a 3-unit and both numerator coefficients have valuation
one. Thus \(v(z(2P))=1\).

\begin{lemma}\label{lem:anchor-index}
If \(R=nP\ne O\) is in the formal kernel at 3, then \(n\)
is even and
\[
             v(z(R))=1+v_3(n),\qquad v(1/x(R))=2+2v_3(n).
\]
\end{lemma}
\begin{proof}
The reduction calculation forces \(2\mid n\). The integral
formal group in parameter \(z=-x/y\) is odd. For a 3-unit
integer \(u\), \([u]z=uz+O(z^3)\) preserves positive
valuation. The expansion \([3]z=3z+O(z^3)\) increases valuation
by exactly one when \(v(z)\ge1\): the remaining terms have
valuation at least \(3v(z)>v(z)+1\). Apply this to
\(R=[n/2](2P)\), also allowing negative indices. Finally the
formal expansion \(1/x=z^2+\text{higher terms}\), with leading
unit coefficient, doubles the valuation.
\end{proof}

For the chain and equations \eqref{eq:anchor}, put
\(m=2^{38347}\). Then \(f_0+f_1\sqrt2=1/x(R)\) and
\(v(1/x(R))\ge2^{38348}=2m\). Thus \(R=nP\) is formal at
3, and Lemma~\ref{lem:anchor-index} gives
\begin{equation}\label{eq:anchor-large}
                       v_3(n)\ge m-1,\qquad |n|\ge3^{m-1}.
\end{equation}
For the cutoff comparison, put \(A=(8D_*)^{288D_*}\). From the
computed bit lengths, \(D_*<2^{38322}\) and
\(\log_2(8D_*)<38325\). Since \(A>6^{100}\) and \(A>128\),
\(H_*<8A(B_*^{\rm exc})^{24}\), and \(L_*<32H_*\). Hence
\begin{equation}\label{eq:anchor-cutoff}
 \log_2L_*<8+288\cdot38325\cdot2^{38322}+24\cdot109445
             <2^{38346}=m/2.
\end{equation}
The last comparison is checked as an integer inequality in the code.
It follows that \(3^{m-1}>L_*\).

Anchor points exist. In each of the finitely many local fields above
2 and 3, the sets of points sufficiently close to \(O\) contain
open subgroups of finite index in the compact group \(E(F_v)\).
Their preimages in \(\Z P\) contain nonzero finite-index
subgroups. Intersect these subgroups, making the point deep enough
at 3 to meet the bound \(2m\), and deep enough at 2 that both
rational coefficients of \(1/x\) are integral. Since \(C\)
is a 2-unit, both \(t_i\) are integral at 2 and 3. Equation
\eqref{eq:T23} supplies the remaining witnesses. This proves
existence without evaluating their enormous coordinates.

\section{Contact forms and multipliers}\label{sec:contacts}

\subsection{Recovering the five contact values}\label{sec:numerical-contact}

Let \(b_m\) be the branch value labeled by
\(m\in\{30,42,70,105,210\}\) in the discriminant-210 quotient.
The four published quartic models for
\(X_{210}/\langle w_m\rangle\), \(m=30,42,70,105\), in
\cite{PadurariuSaia} give the following invariants:
\[
\begin{array}{c|r|c}
m&j(X_{210}/\langle w_m\rangle)&t_m\ \text{with }j_{\rm Leg}(t_m)=j\\\hline
30&191342053882402567201/129708022500&1/2401\\
42&2179252305146449/66177562500&32/375\\
70&128031684631201/9922500&1/225\\
105&1169975873419524361/108425318400&5/1029
\end{array}
\]
Here \(j_{\rm Leg}(t)=256(t^2-t+1)^3/(t^2(1-t)^2)\).
For a quartic \(ax^4+bx^3+cx^2+dx+e\), these values are computed by
\[
 I=12ae-3bd+c^2,\quad
 J=72ace+9bcd-27ad^2-27b^2e-2c^3,\quad
 j=\frac{6912I^3}{4I^3-J^2}.
\]

To justify using these invariants, the residual group
\(\mathcal W/\langle w_m\rangle\cong(\Z/2\Z)^3\) acts
faithfully on the genus-one quotient. Over \(\overline\Q\), its
translation subgroup has order 4 and is the full two-torsion group
of the Jacobian. Quotienting by these translations preserves the
\(j\)-invariant, since it is the multiplication-by-2 isogeny after
choosing an origin. The remaining double cover of \(X^*\) branches
at exactly the four values other than \(b_m\): the inertia at
\(b_m\) was killed by the first quotient. Thus its \(j\)-invariant
is the Legendre invariant of those four values.

Normalize \(b_{210}=\infty\), \(b_{105}=0\), and \(b_{70}=1\).
Set \(\lambda=b_{42}\), \(\mu=b_{30}\). If
\[
 \mathcal O(t)=\{t,1-t,1/t,1/(1-t),t/(t-1),(t-1)/t\},
\]
the first two rows imply
\(\lambda\in\mathcal O(1/2401)\) and
\(\mu\in\mathcal O(32/375)\). The six distinct elements exhaust
the roots of the corresponding degree-six equation for \(j\).
Among the 36 pairs, the equations
\[
 j_{\rm Leg}(\mu/\lambda)=j_{\rm Leg}(1/225),\qquad
 j_{\rm Leg}((\mu-1)/(\lambda-1))=j_{\rm Leg}(5/1029)
\]
have exactly the solution
\[
                    \lambda=2401,\qquad\mu=-343/32.
\]
This is a finite exact rational calculation, reproduced in the
accompanying verification script. It proves the claimed normalization
relative to the published quotient models. In particular we may fix
\begin{equation}\label{eq:numerical-forms}
L_\infty=V,\quad L_0=U,\quad L_1=U-V,\quad
L_{2401}=U-2401V,\quad L_{-343/32}=32U+343V.
\end{equation}
All coefficient and nonzero pairwise-determinant prime divisors are
in \(\{2,3,5,7\}\). The six determinants between finite forms,
up to sign, are \(1,2401,343,2400,375,77175\).

\subsection{A catalogue with 143 small generating primes}

Put \(S_0=\{2,3,5,7,11,13,17,19,23,29\}\),
\(H_0=4\), \(J_0=\prod_{p\in S_0}p=6469693230\), and
\(W=4J_0=25878772920\). The Chinese remainder theorem gives
\[
 (\Z/W\Z)^\times\simeq
 C_2^3\times C_4\times C_6\times C_{10}\times C_{12}
                  \times C_{16}\times C_{18}\times C_{22}\times C_{28}.
\]
Its order is \(4087480320\) and its exponent is \(55440<2^{16}\).
Use the following thirteen disjoint pools \(\mathcal P_\nu\).
\begin{center}\small
\begin{tabular}{r l}
\toprule
\(\nu\)& Primes in \(\mathcal P_\nu\)\\
\midrule
1&31, 37, 41, 43, 47, 53, 59, 67, 71, 73, 83\\
2&89, 97, 101, 103, 107, 109, 113, 127, 131, 137, 139\\
3&149, 151, 157, 163, 167, 173, 179, 181, 191, 193, 197\\
4&199, 211, 223, 227, 229, 233, 239, 241, 251, 263, 271\\
5&277, 281, 283, 293, 307, 311, 313, 317, 331, 337, 347\\
6&349, 353, 359, 367, 373, 379, 383, 389, 397, 401, 419\\
7&421, 431, 433, 439, 443, 449, 457, 461, 463, 467, 479\\
8&487, 491, 499, 503, 509, 521, 523, 541, 547, 557, 563\\
9&569, 571, 577, 587, 593, 599, 601, 607, 619, 631, 643\\
10&647, 653, 659, 661, 673, 677, 683, 691, 701, 709, 719\\
11&727, 733, 739, 743, 751, 757, 761, 769, 773, 797, 811\\
12&821, 823, 827, 829, 839, 853, 857, 859, 863, 877, 881\\
13&883, 887, 907, 911, 919, 929, 941, 947, 967, 971, 983\\
\bottomrule
\end{tabular}
\end{center}

Each pool generates \((\Z/W\Z)^\times\). Here is a short,
reproducible certificate instead of enumerating its four billion
elements. Take coordinates for \((\Z/8\Z)^\times\) in the
basis \((-1,5)\), and discrete-log coordinates at every odd prime
of \(S_0\), using its least positive primitive root. For each
\(\ell\in\{2,3,5,7,11\}\), retain the coordinates whose cyclic
orders are divisible by \(\ell\) and reduce them modulo \(\ell\).
In each pool the resulting matrices have ranks
\[
               11,\quad3,\quad1,\quad1,\quad1,
\]
respectively. These are the full dimensions of the quotients by
\(\ell\). The ranks and every coordinate are recorded in
\texttt{generator\_certificate.json}. Gaussian elimination over the
five finite fields verifies them using exact arithmetic. If a subgroup
of a finite abelian group surjects onto each of these quotients, it is
the whole group: a nontrivial quotient would have a quotient of prime
order, contradicting the corresponding surjectivity.

Write the eleven cyclic orders as \(d_j\). Their least common
multiple is 55440. For every
pool \(\nu\) and every cyclic basis vector \(e_j\), solve
\[
 \sum_{q\in\mathcal P_\nu}e_{\nu jq}\,\operatorname{coord}(q)
       =e_j\quad\hbox{in }\prod_j\Z/d_j\Z,
       \qquad 0\le e_{\nu jq}<55440.
\]
The program solves the linear systems modulo \(2,3,5,7,11\), lifts
through their powers, and combines them by CRT. It checks every one of
the 143 resulting residue identities exactly modulo \(W\). Define
\[
             g_{\nu j}=\prod_{q\in\mathcal P_\nu}q^{e_{\nu jq}}.
\]
These are fixed positive integers represented by multiplication gates;
their expanded values are not used as coefficients.

At a contact of type \(\nu\), use
\begin{equation}\label{eq:cyclic-multiplier}
 k=\prod_{p\in S_0}p^{a_p}\prod_{j=1}^{11}g_{\nu j}^{b_j},
 \qquad 0\le a_p\le3,\qquad
 0\le b_j<2^{\lceil\log_2d_j\rceil}.
\end{equation}
This requires 20 bits for the ten small-prime exponents and 35 for
the cyclic exponents, or 55 in total. Each interval for \(b_j\)
contains \(0,\ldots,d_j-1\), so every required unit residue is
represented. The extra choices are harmless: all their prime factors
are in the same prescribed pool. The following proposition proves that these multipliers suffice.

The circuit shares the fixed prime powers and basis products globally.
For a chosen power \(g\), a Boolean bit \(b\) selects the factor
\(1+(g-1)b\). A further sign bit implements the signed raw contact.
The 55 count excludes the sign bit and all gate variables.

\begin{proposition}[Composite multipliers suffice]
The catalogue \eqref{eq:cyclic-multiplier} supplies the finite multipliers
in \textbf{(I4)}, using the source counting and reduction-order results. The exceptional
set can then be taken to be
\begin{equation}\label{eq:small-S}
                         S_* =\{p:p\le76176\}.
\end{equation}
\end{proposition}
\begin{proof}
For each integer \(A\), use coordinates \(T=(u_1,u_2,v)\)
and substitute \(u_3=Av-u_1-u_2\). The thirteen contact forms
\(F_\nu(T)\) are primitive and pairwise independent modulo every
prime greater than 29: their coefficient vectors are \((0,0,1)\)
and \((\alpha,0,\gamma),(0,\alpha,\gamma),
(-\alpha,-\alpha,\alpha A+\gamma)\), with the determinants
computed above. At \(p\in S_0\), the union bound
\(13p^{-4}<1\) gives a class modulo \(p^4\) on which all
\(v_p(F_\nu)<4\). Fix a lift modulo \(p^7\).
Choose a real box avoiding the zero planes, fixing signs
\(\epsilon_\nu\). Remove the resulting \(S_0\)-parts
\(k_{0\nu}=\prod p^{v_p(F_\nu)}\).

As in \cite[Section 6.2]{H10}, the desired positive prime residues
are 1 modulo 8, 3 modulo 7, and 1 modulo the other primes in \(S_0\).
The chosen lifts determine the signed remaining unit classes. Since each
pool generates \((\Z/W\Z)^\times\), select its product \(h_\nu\)
so division by \(k_{0\nu}h_\nu\) gives these residues. The cyclic basis products ensure that the selections lie in
\eqref{eq:cyclic-multiplier}. It remains to impose congruences at the
pool primes without losing local admissibility.

Write \(k_\nu=k_{0\nu}h_\nu\) for the selected multiplier.
If \(q^a\Vert h_\nu\), only this multiplier is divisible by
\(q\), because the pools are disjoint. On the plane
\(F_\nu(T)=0\pmod q\), exclude
\(\epsilon_jF_j(T)=0,-k_j\) for every \(j\ne\nu\).
Each exclusion removes at most a line, and \(24q<q^2\) since
\(q>29\). Lift a remaining point to a solution of
\(F_\nu(T)=0\pmod{q^a}\). Primitivity permits a further lift
for which
\[
       \epsilon_\nu F_\nu(T)/q^a
                   \not\equiv0,-k_\nu/q^a\pmod q.
\]
Thus the divided \(\nu\)-th form is also neither 0 nor \(-1\)
modulo \(q\). The Chinese remainder theorem combines all choices.
Use a modulus divisible by \(q^{a+1}\), and powers \(p^7\)
at primes of \(S_0\). Then its quotient by any \(k_\nu\)
is divisible by every prime in its support, as required in the
local-density argument. The finitely many unused pool primes may be
included among the other fixed excluded primes.

The divided forms are still integral and pairwise nonproportional.
For any remaining prime \(p>29\), the prime-counting local
factor satisfies
\[
 \beta_p\ge(1-1/p)^{-13}(1-13/p)>0,\qquad
 \log\beta_p\ge-154/p^2.
\]
The logarithmic inequality follows from
\(\log(1-x)\ge-x-x^2/(2(1-x))\). Hence their product has a
positive uniform lower bound; later sieve cutoffs can be enlarged without
changing the fixed catalogue.
The subsequent counting and sieve steps use only these properties,
the established residues, and a fixed finite set of multiplier primes;
they are unchanged. All new multiplier primes are at most 983, hence
already lie in \(S_*\). The other exclusions are the geometric primes \(2,3,5,7\), the
elliptic bad and ramified primes \(2,7\), logarithm denominators
(and primes at most 49), and the parity exceptions at most 76176.
\end{proof}

\section{The numerical contact-height bound}\label{sec:height-certificate}
Set
\[
 g_3=|\mathrm{GL}_4(\F_3)|=24261120,\qquad
 D_*=1152!\,g_3^{1152},\qquad
 B_*^{\rm exc}=\prod_{p\le76176}p,
\]
and define
\begin{equation}\label{eq:optimized-H}
 H_*=2\bigl((8D_*)^{288D_*}+6^{100}\bigr)(B_*^{\rm exc})^{24}+128,
 \qquad L_*=18H_*+12.
\end{equation}
We prove that \(h(s)\le H_*M_{S_*}(s)^{24}\) for
\(s\in\mathbf P^1(\Q)\setminus\mathcal C\), by specifying
the cover, comparing its metrics, and applying the arithmetic height bounds.

\subsection{A specified level and the geometric exceptional primes}
\label{sec:geometric-primes}

Let \(\mathscr X\) be the maximal-level moduli stack for the
quaternion algebra of discriminant \(210\), and let \(X\) be its
coarse curve. Use the maximal order and a polarization element
\(\mu^2=-210\). Such an element can be put in a maximal order by
the quaternion embedding criterion; conjugacy of maximal orders gives
the same Shimura curve. Yuan's integral moduli description
\cite[Section 2.1]{YuanKS} makes \(\mathscr X\) a proper, regular,
semistable Deligne--Mumford stack over \(\Z\), smooth outside
\(210\), carrying a universal principally polarized abelian surface.
The source identifies \(X\) as a genus-five curve without elliptic
stabilizers. In particular its generic moduli stabilizers are exactly
\(\{\pm1\}\): the only possible additional finite-order units
would involve \(\Q(i)\) or \(\Q(\sqrt{-3})\), excluded by the
ramified primes 5 and 7 respectively.

Take principal quaternionic level 3, namely
\[
 U(3)=\ker\bigl(\widehat{\mathcal O}^{\times}
                  \longrightarrow(\mathcal O/3\mathcal O)^{\times}\bigr).
\]
The corresponding fine scheme \(Y\) is proper over \(\Z[1/3]\)
and smooth outside \(210\). At 3 the maximal order is local, with
residue field \(\F_9\); its quotient modulo 3 has 81 elements,
of which 9 are nonunits. Thus
\[
 [\widehat{\mathcal O}^{\times}:U(3)]=72,
 \qquad \deg(Y\longrightarrow X)\le72,
 \qquad \deg(Y\longrightarrow\mathbf P^1)\le1152.
\]
These deliberately loose degree bounds include all components.

Here is a justification that this particular choice introduces no
unlisted bad primes. Work over \(R=\Z[1/210]\). The level-change
group has order 72, invertible in \(R\), and its central subgroup
\(\{\pm1\}\) acts trivially on the fine moduli scheme. After
dividing out this kernel, the action on the generic fiber is free.
It stays free: the fixed locus of a nontrivial tame automorphism of
a smooth relative curve, when nonempty, is smooth of relative
dimension zero and therefore lifts a special fixed point to the
generic fiber. This contradicts generic freeness. Consequently
\(Y_R\to X_R\) is finite \'etale and \(X_R\) is smooth and proper.

The Atkin--Lehner action on \(X\) extends over \(R\), by uniqueness
of the smooth stable model of a curve of genus at least two. Its
order 16 is invertible in \(R\). Stabilizers in any geometric
fiber are cyclic, since a tame stabilizer acts faithfully on the
tangent line; as subgroups of \((\Z/2\Z)^4\), they have order at
most two. Fixed loci lift as above. Distinct fixed loci cannot
collide, since that would create a noncyclic stabilizer; the fixed
points of one involution form an \'etale scheme. It follows that
the quotient is a smooth genus-zero curve over \(R\), with five
disjoint branch sections and ramification index two along each.
The rationality and labels of the generic branch values identify
these sections individually. The three sections labelled
\(\infty,0,1\) trivialize the quotient as \(\mathbf P^1_R\).
The other two are the values computed in
Appendix~\ref{sec:numerical-contact}.

After an unramified extension at a prime outside \(210\), the
cover is therefore either \'etale or has local equation
\(w=u^2\) at a branch section, up to a unit. Pullback along a
rational base point of even contact has no ramification. The
universal abelian scheme on \(Y_R\) also extends at every resulting
integral point. This proves that the source fiber-field and
good-reduction arguments work with the specified choices
\begin{equation}\label{eq:geometric-data}
                 e=1152,\qquad S_{\rm geo}=\{2,3,5,7\}.
\end{equation}
This assertion uses the integral moduli theorem and the quotient
identification, not merely the prime divisors of the contact
determinants.

\subsection{Two integral bicanonical sections}
\label{sec:bicanonical}

Put \(t=32s\), and write
\[
 (a_1,a_2,a_3,a_4)=(0,32,76832,-343),\qquad
 P_4(t)=\prod_{i=1}^4(t-a_i).
\]
The functions below mean their pullbacks to \(X\) or to
\(\mathscr X\). Define
\[
             q_0=\frac{(dt)^2}{P_4(t)},\qquad q_1=tq_0.
\]
Both are holomorphic bicanonical sections on the generic fiber
and have no common zero. At a finite branch value one can use
\(t-a_i=u^2\): the zero of \((dt)^2\) cancels that of
\(P_4(t)\). At infinity use \(t=u^{-2}\); then \(q_0\)
has a zero of order two and \(q_1\) is nonvanishing. Away from
these points the map is unramified. In particular \(q_1/q_0=t\).

Let \(\omega=\omega_{\mathscr X/\Z}\), with its Petersson metric.
For each rational prime put
\[
 \Delta_p=\max_{i<j}v_p(a_i-a_j),\qquad
 J=\prod_p p^{\Delta_p}=2^{10}3^2 5^3 7^4=2765952000.
\]

\begin{lemma}\label{lem:integral-bicanonical}
The sections \(J^2q_0,J^2q_1\) extend to global sections of
\(\omega^{\otimes2}\) on \(\mathscr X\).
\end{lemma}
\begin{proof}
Regularity reduces extension to codimension-one points on \'etale
charts. Horizontal divisors were handled above. At the generic
point of a vertical component over \(p\), semistability makes
\(p\) a uniformizer and the relative differential sheaf equals
\(\omega\). For every rational function \(u\), writing
\(u=p^kv\) with \(v\) a unit gives \(v(du)\ge v(u)\),
since \(dp=0\). Consequently, with \(r_i=v(t-a_i)\),
\(v(dt)\ge\max_i r_i\).

If \(v(t)=k<0\), all four \(r_i=k\), so
\(v(q_0)\ge-2k\) and \(v(q_1)\ge-k\), both nonnegative.
Otherwise all \(r_i\ge0\); at most one exceeds \(\Delta_p\).
Writing \(m=\max r_i\), if \(m\le\Delta_p\), then
\(2m-\sum r_i\ge-2\Delta_p\). If \(m>\Delta_p\), then
\(\sum r_i\le m+3\Delta_p\), which gives the same bound.
Thus \(v(q_0),v(q_1)\ge-2\Delta_p\) in all cases, proving
integrality after multiplication by \(J^2\).
\end{proof}

\subsection{An explicit archimedean metric bound}
\label{sec:metric-bound}

Set \(A_*=76832\) and
\[
 C_{\rm met}=16(12A_*+9)^3
             =12540153541944538512.
\]
For \(t\in\C\setminus\{a_1,\ldots,a_4\}\), put
\(r_i=|t-a_i|\), \(U=1+\sum_i r_i\), and use the conformal
metric
\[
                    \rho^2|dt|^2=\frac{U}{\prod_i r_i}|dt|^2.
\]
Its pullback to \(X(\C)\) is smooth and positive, including
the branch points: in a coordinate \(t-a_i=u^2\) the apparent
singularity cancels, and at infinity \(t=u^{-2}\) gives the
coefficient
\[
      4\,\frac{|u|^2+\sum_i|1-a_i u^2|}
                    {\prod_i|1-a_i u^2|}.
\]

Let \(K_\rho\) denote its Gaussian curvature. Off the branch
points the Euclidean Laplacian gives
\[
 -K_\rho=\frac{\prod_i r_i}{2U}\,\Delta\log U
 \ \ge\ \frac{\sum_i\prod_{j\ne i}r_j}{2U^3}.
\]
Indeed \(\Delta U=\sum_i r_i^{-1}\), \(|\nabla U|^2\le16\),
and \((\sum r_i)(\sum r_i^{-1})\ge16\).
For \(|t|\le2A_*+2\), we have \(U\le12A_*+9\).
The branch values are pairwise at distance at least one, so at
most one \(r_i\) is less than \(1/2\). One triple product is
therefore at least \(1/8\), and
\(-K_\rho\ge C_{\rm met}^{-1}\).
For \(r=|t|>2A_*+2\), we have \(r/2\le r_i\le3r/2\)
and \(U\le7r\). Thus \(-K_\rho\ge1/1372\), a stronger bound.
The bound extends to the branch points by continuity.

The Ahlfors--Schwarz lemma, or the maximum principle applied on
the compact curve, now yields
\(\rho^2/C_{\rm met}\le\rho_{\rm hyp}^2\), where the hyperbolic
metric has curvature \(-1\). Yuan's Petersson convention is
\(\|d\tau\|_{\rm Pet}=2\operatorname{Im}\tau\) on the upper
half plane. Consequently
\begin{equation}\label{eq:bicanonical-sup}
 \|q_0\|_{\rm Pet}\le\frac{4C_{\rm met}}U\le4C_{\rm met},
 \qquad
 \|q_1\|_{\rm Pet}\le\frac{4C_{\rm met}|t|}U\le4C_{\rm met}.
\end{equation}
Here \(a_1=0\) ensures \(U\ge1+|t|\). These are norms of
the bicanonical sections themselves, not their squared norms.

\subsection{Comparison with Faltings height}

\begin{proposition}\label{prop:numerical-geometric-height}
For a geometric point \(y\) above a rational base point \(s\),
and the corresponding quaternionic abelian surface \(A_y\),
\[
                 h(s)\le2\max\{0,h_{\rm F}(A_y)\}+128.
\]

\end{proposition}
\begin{proof}
Yuan's Theorem 1.1 is the isometry
\[
  \omega_{A}^{\otimes2}\simeq210\,\omega^{\otimes2}
\]
of hermitian line bundles on the maximal stack. His Hodge metric
has squared norm \((2\pi)^{-2}|\int\alpha\wedge\bar\alpha|\).
Its associated height is
\(h_{\rm F}^{Y}=h_{\rm F}+\log\pi\) in dimension two, where
\(h_{\rm F}\) has the original Faltings normalization used in
\cite{vonKaenelKret}; compare their equation (11.28).
In particular
\[
              h_{\bar\omega^{\otimes2}}(y)
                  =2h_{\rm F}(A_y)+2\log\pi+\log210.
\]
The sign of the last term follows because the sublattice
\(210\omega^{\otimes2}\), with its induced metric, has height
\(h_{\bar\omega^{\otimes2}}-\log210\).
To evaluate these heights, extend the point to the proper maximal
stack after a finite field extension. The universal abelian scheme
then computes stable Faltings height; normalized heights are unchanged
by that extension.

The integral sections in Lemma~\ref{lem:integral-bicanonical}
have no common zero in characteristic zero and ratio \(t\).
At a finite place their common base ideal can only decrease the
height of this ratio. At an infinite place
\eqref{eq:bicanonical-sup} bounds both section norms. The usual
local expression for the height of a projective ratio therefore gives
\[
 h(t)\le h_{\bar\omega^{\otimes2}}(y)
                         +2\log J+\log(4C_{\rm met}).
\]
Since \(s=t/32\), we conclude that
\[
 h(s)\le2h_{\rm F}(A_y)
       +\log210+2\log J+\log(4C_{\rm met})+\log32+2\log\pi.
\]
For an exact bound on the additive term, use \(\pi<4\) and
the integer inequality
\[
             210J^2(4C_{\rm met})\,32\cdot16<2^{160}.
\]
The first five terms in the Taylor expansion of \(e^3\) sum to
\(131/8>16\), so \(\log2<3/4\). The additive term is thus
less than \(120<128\). Replacing \(h_{\rm F}\) by its positive
part proves the assertion.
\end{proof}

\subsection{The field of homomorphisms}
\begin{proposition}\label{prop:smaller-degree}
For the fiber above \(s\), there is a Galois field \(L/\Q\)
of degree at most \(D_*\), defining all geometric homomorphisms
between its conjugate surfaces. It is unramified outside
\(S_*\) and the odd contact primes of \(s\).
\end{proposition}
\begin{proof}
Put \(e=1152\). The fiber has at most \(e\) geometric points. Its splitting field
\(K_0\) has degree at most \(e!\), and there are at most \(e\)
conjugate surfaces \(A_1,\ldots,A_t\), each defined over \(K_0\).
Adjoin the full 3-torsion of each. The resulting field \(L'\) is
Galois over \(\Q\), since the surfaces are permuted by Galois,
and \([L':\Q]\le e!g_3^e\).

All geometric homomorphisms \(A_i\to A_j\) are defined over
\(L'\). Indeed, if a homomorphism kills \(A_i[3]\), it factors
through \([3]\), so is three times another homomorphism. Thus
reduction on 3-torsion embeds
\(\operatorname{Hom}(A_i,A_j)/3\) into
\(\operatorname{Hom}(A_i[3],A_j[3])\). The Galois action on
the integral Hom lattice has finite image. Over \(L'\) it is
congruent to the identity modulo 3, hence trivial by torsion-freeness
of the principal congruence kernel at 3. This is the usual torsion
argument for fields of definition of homomorphisms
\cite{SilverbergHom}.

At good primes of residue characteristic different from 3, the
3-torsion extension is unramified. The preceding model calculation makes the fiber field unramified
outside \(S_*\) and its odd contact primes. Since 3 is exceptional,
adjoining 3-torsion preserves this support.
\end{proof}

\subsection{The arithmetic height estimate}
\begin{proof}[Proof of the height bound \eqref{eq:height}]
Put \(M=M_{S_*}(s)\), \(B_S=B_*^{\rm exc}\), and \(D_{\rm hom}=D_*\). The splitting field of the fiber, followed by adjoining the full
3-torsion of its conjugate surfaces, has degree at most
\(D_{\rm hom}\). Proposition~\ref{prop:smaller-degree} proves that
this field, denoted by \(L\), is Galois and defines every geometric
homomorphism between these surfaces. At good primes outside \(S_*\), the torsion extension
is unramified; the source's odd-contact argument controls the fiber
splitting field. Thus
\[
 \operatorname{rad}(D_L N_{\rm bad})\mid B_S M,
\]
where \(D_L\) is the absolute field discriminant and
\(N_{\rm bad}\) is the product of the norms of the bad-reduction
primes of the surface over \(L\).

For a non-CM fiber its endomorphism algebra is the rational
quaternion algebra. Skolem--Noether adjusts each conjugate isogeny
to commute with this algebra. A quadratic subfield then supplies
the required \(\mathrm{GL}_2\)-structure with compatible
\(G_\Q\)-isogenies. All these maps are defined over \(L\),
so the normality and field-of-definition hypotheses hold.
Proposition 7.13 of \cite{vonKaenelKret}, in dimension 2, gives
\[
 h_{\rm F}(\mathcal A_y)
       \le(8D_{\rm hom})^{288D_{\rm hom}}(B_SM)^{24}.
\]
For a CM fiber we keep the field degree in the isogeny estimate.
Put \(R=B_SM\) and \(d_L=[L:\Q]\le D_{\rm hom}\).
The maximal-order CM factors in the proof of
\cite[Section 7.3.3]{vonKaenelKret}, using Lemmas 7.10 and 7.12
there, give an \(L\)-isogenous surface \(A'\) with
\[
                       h_{\rm F}(A')\le V:=2\cdot6^{96}R^{10}.
\]
In more detail, each simple maximal-order factor has height at most
\(6^{96}\operatorname{rad}(D_L)^{10}\), and at most two factors,
counted with multiplicity, occur. All can be taken over \(L\)
because the endomorphisms are defined there.
The isogeny-height estimate \cite[(7.13)]{vonKaenelKret} then gives
\[
 h_{\rm F}(\mathcal A_y)
 \le V+8\log\bigl(14^{32}d_L\max(V,\log d_L,1)\bigr).
\]
Since \(\max(V,\log d_L,1)\le V D_{\rm hom}\),
\(\log14<3\), \(\log6<2\), and \(\log R\le R^{10}\),
the right side is at most
\[
 (2\cdot6^{96}+2392+16D_{\rm hom})R^{10}
     \le(6^{100}+16D_{\rm hom})R^{10}.
\]
This retains the degree-dependent cost of passing from maximal CM
orders to the actual surface. Finally
\((8D_{\rm hom})^{288D_{\rm hom}}\ge16D_{\rm hom}\).
Thus both cases are bounded by
\(\bigl((8D_*)^{288D_*}+6^{100}\bigr)(B_SM)^{24}\).
Proposition~\ref{prop:numerical-geometric-height} and \(M\ge1\)
give the required inequality \eqref{eq:height}.
\end{proof}

Exact arithmetic computes
\[
 \operatorname{bitlength}(D_*)=38322,\qquad
 \operatorname{bitlength}(B_*^{\rm exc})=109445,
 \qquad |S_*|=7499.
\]
Both integers, in full hexadecimal form, accompany the code in
\texttt{compact\_constants.json}. Neither \(H_*\) nor \(L_*\)
has been numerically expanded. 

\section{Parity and fixed local rings}
\label{sec:parity}
\subsection{The parity polynomial}

For this section and the final integer definition, use the simultaneous
point-realization result supplied by \cite[Lemmas 3.4 and 3.6]{H10},
restated in \cite[Arithmetic Input 2.3]{Conant}. Explicitly, for positive squarefree \(\beta,\theta\),
their supports \(B,J\) are disjoint and avoid \(2,3\), with
\(B\ne\varnothing\); \(\theta\) is a local square at every
prime of \(B\), and \(\beta\) at every prime of \(J\).
Both are squares in \(\Q_2\), while \(\theta\) is a square
in \(\Q_3\) and \(\beta\) has the unit nonsquare class there.
These results supply a positive squarefree \(n\) making both
\[
 E_n:y^2=x(x-n)(x+3n),\qquad E_{n\theta}
\]
have rational points with nonzero \(x\)-coordinate in
\(\beta\Q^{\times2}\). The different notation \(E_l\) denotes
this auxiliary family, not the fixed curve over \(F\).

Define a completely numerical polynomial
\[
                  A_N(T)=\prod_{j=1}^N(T^2-j^2).
\]
For the rational input \(b\), take eleven rational witnesses
\[
                  \mu,\delta,t,U,V,e,s_1,s_2,y_1,y_2,w.
\]
Abbreviate \(b'=\mu b\) and \(X_i=b's_i^2\), and put
\begin{align*}
 f_1&=A_{138}(\mu),&f_2&=A_{76176}(\delta),\\
 f_3&=(1+\delta bt^2)U-b(1-\delta bt^2),\\
 f_4&=V-b'(b'-U^2),\\
 f_5&=y_1^2-X_1(X_1-e)(X_1+3e),\\
 f_6&=y_2^2-X_2(X_2-eV)(X_2+3eV),\\
 f_7&=(1+\delta bt^2)eV\mu b s_1s_2w-1.
\end{align*}
Set
\begin{equation}\label{eq:numerical-phi}
 P_\Phi(b,\mathbf z)=b\sum_{i=1}^7f_i^2,
 \qquad\Phi(b):\ \exists\mathbf z\in\Q^{11}\ P_\Phi(b,\mathbf z)=0.
\end{equation}
The abbreviations are polynomial substitutions. Thus this is one
specified polynomial with integer coefficients, of total degree
\(304705\). The product expression specifies all its coefficients
without expanding them. It is an auxiliary parity predicate,
\emph{not by itself a definition of the integers}.

\begin{proposition}[Numerical normalization]\label{prop:explicit-phi}
With \(D,S_*\) as in Section~\ref{sec:inputs},
\[
                         D\subseteq\Phi(\Q)\subseteq D_{S_*}.
\]
\end{proposition}

\begin{proof}
For soundness, let \(q>76176\), and suppose \(k=v_q(b)<0\) is
odd. The two integer multipliers \(\mu,\delta\) are \(q\)-units.
The valuation of \(\delta bt^2\), when \(t\ne0\), is odd and
nonzero, so \((1-\delta bt^2)/(1+\delta bt^2)\) is a unit.
Consequently \(v_q(U)=k\) and \(v_q(V)=3k\). Exactly one of
\(e,eV\) has even valuation. After scaling that curve to a unit
parameter at \(q\), every nonzero \(x\)-coordinate has even
valuation: a negative valuation gives \(v_q(y^2)=3v_q(x)\), and
a positive one gives \(v_q(y^2)=v_q(x)\). This contradicts the
odd valuation of its required square class \(b'\Q^{\times2}\).

For completeness, first justify the numerical multiplier bounds.
For each prescribed sign, the numbers
\[
 \mu=\pm2^{e_2}3^{e_3}u,\qquad e_2,e_3\in\{0,1\},\quad
 u\in\{1,5,7,11,13,17,19,23\},
\]
represent all 32 pairs of local square classes at \(2\) and \(3\).
Indeed, the eight values of \(u\) are all the units modulo 24.
Every \(|\mu|\le138\). Choose one with \(b'=\mu b>0\),
square class 1 at \(2\), and the unit nonsquare class at \(3\).
Its prime support, besides \(2,3\), contains at most one prime
\(p\in\{5,7,11,13,17,19,23\}\).

For this set \(S(\mu)\), representatives of all local square-class
tuples with a prescribed real sign may be chosen as
\(\delta=\pm wu\), where \(w\) is a squarefree product of
primes in \(S(\mu)\), and \(u\) is a positive unit residue
modulo \(8\prod_{q\in S(\mu),\ q\ne2}q\). For each odd prime
choose either a square or nonsquare unit residue, and use the Chinese
remainder theorem; at \(2\), use the four unit residues modulo 8.
When \(p\) occurs, \(w\le6p\) and \(u<24p\), so
\[
                         |\delta|<144p^2\le76176.
\]
Without \(p\) the smaller bound 138 suffices. Thus choose an
allowed \(\delta\) with \(\delta b>0\) and a square in
\(\Q_q\) for every \(q\in S(\mu)\).

Choose \(t\) by weak approximation, close to a solution of
\(\delta bt^2=1\) at the real place and at every prime of
\(S(\mu)\). Then \(U\) is sufficiently small that
\(V=b'^2(1-U^2/b')\) is positive and is a local square at those
primes. At every prime of \(\Supp(\delta)\setminus S(\mu)\),
add the condition \(v_q(\delta bt^2)>0\). Avoid the finitely many
zero denominators and \(V=0\).

Let \(\beta,\theta\) be the positive squarefree representatives
of \(b',V\). If \(q\in\Supp(\beta)\setminus S(\mu)\),
the valuation \(v_q(b)=v_q(b')\) is odd and strictly positive,
because \(b\in D\). The quotient defining \(U/b\) is a unit,
either by parity or by the extra smallness condition. Thus
\(v_q(U^2/b')>0\), and \(V\) is a local square at \(q\).
This also shows that the two squarefree supports are disjoint.

Conversely, if \(q\in\Supp(\theta)\), then \(q\notin S(\mu)\)
and \(v_q(b')\) is even. Odd valuation of \(V\) forces
\(v_q(b')=2v_q(U)\) and \(v_q(b'-U^2)>v_q(b')\).
Hence \(b'/U^2\in1+q\Z_q\), so \(b'\) is a local square.
If \(U=0\), the support of \(\theta\) is empty and there is
nothing to check. The conditions at \(2,3\) and positivity are
already arranged, and \(\Supp(\beta)\ne\varnothing\) because
of its nonsquare class at \(3\).

The cited point-realization result now gives points on \(E_n,E_{n\theta}\)
in square class \(b'\). Write \(V=\theta c^2\); the change
\((x,y)\mapsto(c^2x,c^3y)\) gives the second point on
\(E_{nV}\). Take \(e=n\), choose \(s_i,y_i\) from these points,
and let \(w\) invert the nonzero product in \(f_7\). All seven
equations vanish. For \(b=0\), \eqref{eq:numerical-phi} holds
automatically.
\end{proof}

\subsection{The fixed local rings}\label{sec:local-rings}

For fixed nonzero rational \(a,b\), define
\[
 S_{a,b}=\{2z_1:\exists z_2,z_3,z_4\in\Q,
             \ z_1^2-az_2^2-bz_3^2+abz_4^2=1\},
 \qquad T_{a,b}=S_{a,b}+S_{a,b}.
\]
Koenigsmann's Proposition 10 \cite{Koenigsmann} gives the following
choices for the set defined by \(\Int_p\):
\[
\begin{array}{c|l}
 p&\Z_{(p)}\\\hline
 2&T_{3,3}+T_{2,5}\\
 p\equiv3\pmod8&T_{-1,-p}+T_{2,-p}\\
 p\equiv5\pmod8&T_{-2,-p}+T_{2,-p}\\
 p\equiv7\pmod8&T_{-1,-p}+T_{-2,p}\\
 p\equiv1\pmod8&T_{-2p,q}+T_{2p,q}
\end{array}
\]
In the last row, choose the least prime \(q\equiv3\pmod8\) with
Legendre symbol \((p/q)=-1\). Such primes exist, and trial search
determines this fixed coefficient. Sums mean existentially quantified
summands. Each row therefore expands directly into polynomial equations.
These definitions concern only the finitely many fixed primes in \(S\).

For a rational \(t\), let \(T_{2,3}(t)\) be the seven-witness
predicate
\begin{equation}\label{eq:T23}
\begin{split}
 u_0^2-2u_1^2-3u_2^2+6u_3^2&=1,\\
 (t-2u_0)^2-8v_1^2-12v_2^2+24v_3^2&=4.
\end{split}
\end{equation}
The quaternion \((2,3)\) ramifies precisely at 2 and 3 and is
split over \(\mathbb R\). Proposition 6 of \cite{Koenigsmann}
therefore gives
\(T_{2,3}(\Q)=\Z_{(2)}\cap\Z_{(3)}\).


\begin{thebibliography}{99}\footnotesize\setlength{\itemsep}{2pt}
\bibitem{H10} OpenAI, \emph{Hilbert's tenth problem over the rational
numbers}, preprint, version of 6 October 2026 (original version
24 September 2026).
\href{https://github.com/openai/math/blob/main/preprints/Hilberts-tenth-problem-over-the-rational-numbers-October-6-2026/main.pdf}{Source manuscript}.
Propositions 2.1, 2.2 and 2.8; Theorem 2.4; Corollaries 2.6 and 6.2;
Lemmas 3.4, 3.6 and 6.1. 
\bibitem{Conant} \emph{Existential definability of the integers from
elliptic indices and height bounds}, AI-generated manuscript supplied
by Gabe Conant, 7 October 2026. In particular, Arithmetic Input 2.3,
Proposition 3.1, and Sections 5--7.
\bibitem{ConantArity} \emph{Simultaneous local certificates and a
reduced existential arity}, AI-generated research draft supplied by
Gabe Conant, 8 October 2026, eight pages. Lemma 2.1 and the shared-coordinate construction are used
in Section~\ref{sec:simultaneous}.
\bibitem{Koenigsmann} J. Koenigsmann, \emph{Defining $\Z$ in $\Q$},
Annals of Mathematics \textbf{183} (2016), 73--93.
\href{https://doi.org/10.4007/annals.2016.183.1.2}{DOI}.
\bibitem{PadurariuSaia} O. Padurariu and F. Saia,
\emph{Shimura curve Atkin--Lehner quotients of genus at most two},
arXiv:2509.25368, 2025. The four discriminant-210 quartics used here are
in their accompanying data file
\href{https://github.com/fsaia/GenusAtMost2/blob/main/non_elliptic_genus_1_equation_determined_final.m}{published model data}.
Their identification as quotient models is used as a published input.
\bibitem{vonKaenelKret} R. von K\"anel and A. Kret,
\emph{Integral points on coarse Hilbert moduli schemes},
arXiv:2307.06944, 2023, Proposition 7.13, Lemmas 7.10 and 7.12,
equation (7.13), and the maximal-order estimates in Section 7.3.3.
\href{https://arxiv.org/abs/2307.06944}{Preprint}.
\bibitem{YuanKS} X. Yuan,
\emph{Explicit Kodaira--Spencer map over Shimura Curves},
arXiv:2205.11334v2, 2024, Theorem 1.1 and Section 2.
\href{https://arxiv.org/abs/2205.11334}{Preprint}.
\bibitem{DPRMCoq} D. Larchey-Wendling and Y. Forster,
\emph{Hilbert's Tenth Problem in Coq (Extended Version)},
Logical Methods in Computer Science \textbf{18}(1), 2022,
article 35. \href{https://arxiv.org/abs/2003.04604}{Preprint}.
\bibitem{SilverbergHom} A. Silverberg,
\emph{Fields of definition for homomorphisms of abelian varieties},
Journal of Pure and Applied Algebra \textbf{77} (1992), no. 3,
253--262. \href{https://www.sciencedirect.com/science/article/pii/0022404992901412}{Publisher record}.
The torsion-field inclusion is also proved in the form needed here
in Proposition~\ref{prop:smaller-degree}.
\bibitem{DPRMLibrary} The Library of Undecidability,
\href{https://github.com/uds-psl/coq-library-undecidability/tree/8880e198bdc44cba1bd901f1fee701a95fc440ae}{pinned \texttt{coq-8.18} source revision},
used for the compilation of 8 October 2026.
The source files under \texttt{theories/H10} used here are
\texttt{Matija/cipher.v}, \texttt{Dio/dio\_cipher.v},
\texttt{Dio/dio\_bounded.v}, \texttt{Dio/dio\_expo.v}, and
\texttt{Dio/dio\_elem.v}. The accompanying package contains the extracted constraints, wrapper
proofs, reconstruction maps and execution records.
\end{thebibliography}
\end{document}
